\documentclass[11pt,a4paper]{article}

\usepackage[margin=2.5cm]{geometry}
\usepackage{enumitem}
\usepackage{titlesec}
\usepackage{parskip}
\usepackage{xcolor}
\usepackage{hyperref}
\usepackage{fancyhdr}
\usepackage{booktabs}
\usepackage{array}

\hypersetup{colorlinks=false}

% Header / Footer
\pagestyle{fancy}
\fancyhf{}
\rhead{\small Nour Ebid --- Seminar zu aktuellen Entwicklungen}
\lhead{\small HTW Berlin, SS 2026}
\rfoot{\small \thepage}

% Section style
\titleformat{\section}{\large\bfseries}{}{0em}{}[\titlerule]
\titleformat{\subsection}{\normalsize\bfseries\color{black!70}}{}{0em}{}

% List style
\setlist[itemize,1]{leftmargin=1.2em, itemsep=4pt, parsep=0pt}
\setlist[itemize,2]{leftmargin=1.5em, itemsep=2pt, parsep=0pt}

\begin{document}

% ── Title block ───────────────────────────────────────────────────────────────
\begin{center}
    {\LARGE\bfseries What I Did --- Project Progress Summary}\\[6pt]
    {\large Face Authentication: Real vs.\ Synthetic Training Data}\\[10pt]
    \begin{tabular}{rl}
        \textbf{Student:}   & Nour Ebid \\
        \textbf{Program:}   & M.Sc. Angewandte Informatik, Fachbereich 4 \\
        \textbf{Seminar:}   & Seminar zu aktuellen Entwicklungen \\
        \textbf{Betreuer:}  & Frau Prof.\ Knaut \\
        \textbf{University:}& Hochschule f{\"u}r Technik und Wirtschaft Berlin \\
        \textbf{Date:}      & May 2026 \\
    \end{tabular}
\end{center}

\vspace{8pt}
\hrule
\vspace{12pt}

% ── Step 1 ────────────────────────────────────────────────────────────────────
\section*{Step 1 \quad Chose the Research Topic and Set Up the Project}

\begin{itemize}
    \item Decided to explore whether AI-generated (synthetic) face images can replace real photos in a face authentication system
    \item The motivation: real biometric photos raise privacy concerns under GDPR — synthetic images could remove the need to store sensitive data
    \item Set up the project folder, Python virtual environment, and installed the required libraries: \texttt{deepface}, \texttt{insightface}, \texttt{opencv-python}, \texttt{numpy}
    \item Collected \textbf{15 real photos} of each of 4 friends (\textit{amr}, \textit{adel}, \textit{otto}, \textit{mazen}) using a smartphone under normal, everyday conditions
    \item Downloaded \textbf{1,860 AI-generated face images} from \textit{thispersondoesnotexist.com} (StyleGAN2) to use as synthetic base images
\end{itemize}

% ── Step 2 ────────────────────────────────────────────────────────────────────
\section*{Step 2 \quad Built the Base Authentication System}

\begin{itemize}
    \item Built \texttt{auth\_system.py}: a face authentication system using \textbf{FaceNet} accessed via the DeepFace library
    \item How it works:
    \begin{itemize}
        \item \textbf{Registration:} each training photo is passed through FaceNet → produces a 128-dimensional embedding vector → all embeddings for one person are averaged into one reference vector → saved to a JSON file
        \item \textbf{Authentication:} a query photo is embedded → L2 distance is computed against every stored reference → if the nearest distance is below threshold $\theta = 10.0$, access is granted
    \end{itemize}
    \item Chose \textbf{NIST FRTE metrics} (FAR and FRR) as the evaluation framework, since these are the standard for biometric system benchmarking
    \item Defined the research hypothesis: synthetic-data registration will achieve a FAR within 15 percentage points of real-data registration
\end{itemize}

% ── Step 3 ────────────────────────────────────────────────────────────────────
\section*{Step 3 \quad Ran Experiment 1 (Real Photos) as the Baseline}

\begin{itemize}
    \item Registered all 4 friends using their real photos, tested on the same photos, got initial results
    \item Later refined this to a proper train/test split (see Step 7 below)
    \item Final Experiment 1 setup: register on first 10 photos, test on last 5
    \item Results: \textbf{FAR = 10\%, FRR = 0\%, Accuracy = 95\%}
    \begin{itemize}
        \item All 20 genuine test queries (4 friends × 5 photos) correctly verified
        \item 2 out of 20 AI impostors incorrectly accepted (L2 scores: 6.60 and 5.56)
    \end{itemize}
    \item These 2 false acceptances made sense — some AI-generated faces happened to produce embeddings similar to a registered user's cluster
\end{itemize}

% ── Step 4 ────────────────────────────────────────────────────────────────────
\section*{Step 4 \quad First Synthetic Approach: Face Swap (Failed)}

\begin{itemize}
    \item Idea: take my friends' faces and use Poisson blending (\texttt{cv2.seamlessClone}) to paste them onto randomly selected AI base faces → creates a "synthetic-looking" version of each person
    \item Built the pipeline, generated synthetic images, ran Experiment 2 with them
    \item \textbf{Result: FRR = 75\%} — the system rejected 3 out of every 4 genuine friend attempts
    \item Investigated why:
    \begin{itemize}
        \item Found that approximately \textbf{50\% of the downloaded AI base images were female}
        \item All 4 registered friends are male
        \item FaceNet encodes gender-correlated facial geometry (jaw, brow, bone structure) in its embeddings
        \item Blending a male face onto a female AI base produced an embedding that was "halfway between" — too far from the genuine male embedding to pass the threshold
    \end{itemize}
\end{itemize}

% ── Step 5 ────────────────────────────────────────────────────────────────────
\section*{Step 5 \quad Second Attempt: Gender-Filtered Face Swap (Also Failed)}

\begin{itemize}
    \item Thought the fix was simple: filter the AI base pool to only use male-presenting faces
    \item Used InsightFace's gender classification (confidence $> 0.6$) to select only male bases
    \item Re-ran the pipeline and Experiment 2
    \item \textbf{Still failed} — FRR remained high
    \item Diagnosed the new issue:
    \begin{itemize}
        \item The Poisson blending algorithm creates smooth pixel transitions at the face boundary
        \item These transitions introduced \textbf{small low-frequency artifacts} that FaceNet detects
        \item The genuine L2 scores were consistently shifted \textbf{upward by 1--2 units} across all users
        \item With a threshold of 10.0, that shift was enough to push many genuine scores above it → rejections
    \end{itemize}
    \item Conclusion: the face swap approach has a fundamental problem — the foreign base face and the blending process both contaminate the embedding, regardless of which base face is used
\end{itemize}

% ── Step 6 ────────────────────────────────────────────────────────────────────
\section*{Step 6 \quad Redesigned the Synthesis: Augmentation Pipeline}

\begin{itemize}
    \item Abandoned face swapping completely
    \item New approach: instead of replacing the face's background, \textbf{transform the face itself} using controlled imaging parameters
    \item Key insight: lighting, colour temperature, skin tone, camera noise, and pose are all \textit{imaging variables} — they change how a photo looks, but they do not change whose face it is
    \item Built a 5-stage augmentation pipeline (\texttt{generate\_synthetic\_friends.py}):
    \begin{itemize}
        \item \textbf{Stage 1 — Face alignment:} use InsightFace (buffalo\_l model) to detect 5 facial keypoints and warp the face to a standard 256×256 frontal template — removes pose variation
        \item \textbf{Stage 2 — Lighting:} apply gamma correction with $\gamma$ randomly chosen from $[0.55,\, 1.60]$ — simulates bright or dim lighting
        \item \textbf{Stage 3 — Colour temperature:} randomly apply warm, cool, or neutral colour shift — simulates different camera sensors and light sources
        \item \textbf{Stage 4 — Skin tone:} apply CLAHE (adaptive histogram equalisation) to the luminance channel in LAB colour space — normalises skin tone without changing hue
        \item \textbf{Stage 5 — Camera quality and pose:} add Gaussian blur + random noise to simulate camera quality; add small in-plane rotation ($\pm 12°$) to simulate natural head tilt
    \end{itemize}
    \item Generated \textbf{15 synthetic images per person} from their 10 training photos (cycling through with different random parameters each time)
\end{itemize}

% ── Step 7 ────────────────────────────────────────────────────────────────────
\section*{Step 7 \quad Fixed the Train/Test Split (Addressed the 100\% Accuracy Problem)}

\begin{itemize}
    \item An earlier version of the experiment gave \textbf{100\% accuracy on synthetic data} — the professor correctly identified this as unrealistic
    \item The mistake: I had used all 15 photos per person to generate synthetic images, and then tested on those same 15 real photos
    \item This is called \textbf{circular evaluation} — the test images had already been seen during the generation step, so the model had an unfair advantage
    \item Fix: enforced a strict 10/5 split:
    \begin{itemize}
        \item \textbf{First 10 photos} per person → used for registration (Experiment 1) or as source for synthetic generation (Experiment 2)
        \item \textbf{Last 5 photos} per person → held out completely, never touched during either registration or generation
        \item \textbf{Both experiments use the exact same held-out 5 photos} for testing — so any difference in results is due to the training data type only
    \end{itemize}
    \item After this fix, results became realistic and properly comparable
\end{itemize}

% ── Step 8 ────────────────────────────────────────────────────────────────────
\section*{Step 8 \quad Ran Experiment 2 (Synthetic Registration) with the Fixed Setup}

\begin{itemize}
    \item Registered all 4 friends using their 15 synthetic images (generated from their 10 training photos)
    \item Tested on the same held-out 5 real photos per person + same 20 AI impostors as Experiment 1
    \item Results: \textbf{FAR = 0\%, FRR = 5\%, Accuracy = 97.5\%}
    \begin{itemize}
        \item 19 out of 20 genuine queries verified correctly
        \item 1 rejection: one of \textit{otto}'s photos scored 11.47 (just above threshold 10.0)
        \item All 20 AI impostors correctly rejected (FAR = 0\%)
    \end{itemize}
\end{itemize}

% ── Step 9 ────────────────────────────────────────────────────────────────────
\section*{Step 9 \quad Compared the Two Experiments and Evaluated the Hypothesis}

\begin{itemize}
    \item Side-by-side comparison:

    \vspace{6pt}
    \begin{center}
    \begin{tabular}{lccc}
    \toprule
    Metric & Experiment 1 (Real) & Experiment 2 (Synthetic) & Difference \\
    \midrule
    FAR & 10.0\% & 0.0\% & $-$10.0 pp \\
    FRR & 0.0\%  & 5.0\% & $+$5.0 pp \\
    Accuracy & 95.0\% & 97.5\% & $+$2.5 pp \\
    Friends verified & 20/20 & 19/20 & $-$1 \\
    Impostors rejected & 18/20 & 20/20 & $+$2 \\
    \bottomrule
    \end{tabular}
    \end{center}
    \vspace{6pt}

    \item The FAR gap is $-$10 percentage points, which is within the $\pm$15 pp threshold → \textbf{hypothesis confirmed}
    \item Explained the FAR/FRR trade-off: the controlled synthetic pipeline produces tighter embedding clusters → clearer separation from impostors (lower FAR) but slightly less tolerance for natural variation in test photos (slightly higher FRR)
    \item Otto's one false rejection is explained by this: his held-out test photo had more variation than the controlled synthetic training images could fully represent
\end{itemize}

% ── Step 10 ───────────────────────────────────────────────────────────────────
\section*{Step 10 \quad Wrote the Zwischenbericht and Research Paper}

\begin{itemize}
    \item Wrote the full research paper in IEEE conference format (\texttt{IEEEtran}) using LaTeX
    \item Paper includes: Abstract, Introduction, Background and Related Work, Hypothesis and Variables, Methodology (including the rejected approaches), Experimental Results, Discussion, Ethical Considerations, Conclusion and Future Work
    \item Added 12 academic references including FaceNet (Schroff et al., 2015), NIST FRTE, InsightFace/ArcFace, GDPR, GAN papers (Goodfellow, StyleGAN2), and prior synthetic face recognition work (Wood et al., Bae et al.)
    \item Also prepared a separate Zwischenbericht (interim report) marking which sections are complete and which are planned, with a timeline to mid-June
    \item Planned next steps for the final paper:
    \begin{itemize}
        \item Full DET curve evaluation across multiple threshold values
        \item Expanding the dataset (more users, more test images)
        \item Potentially trying a GAN-based approach with fully AI-generated identities
    \end{itemize}
\end{itemize}

\vspace{12pt}
\hrule
\vspace{8pt}
\begin{center}
\small\textit{Nour Ebid --- Seminar zu aktuellen Entwicklungen --- HTW Berlin --- SS 2026}
\end{center}

\end{document}
